\documentclass[11pt,a4paper]{article}

\usepackage[spanish,es-nodecimaldot]{babel}
\usepackage[T1]{fontenc}
\usepackage[utf8]{inputenc}
\usepackage{lmodern}
\usepackage{geometry}
\usepackage{graphicx}
\usepackage{float}
\usepackage{hyperref}
\usepackage{enumitem}

\geometry{margin=2.5cm}
\hypersetup{colorlinks=true,linkcolor=black,urlcolor=blue}

\title{Ingeniería social}
\author{Javier Siñuka}
\date{23 de septiembre de 2026}

\begin{document}

\maketitle

\section{Objetivo}

Investigar qué es la ingeniería social, reconocer algunos de sus métodos y
conocer formas básicas de prevenir estos ataques.

\section{Respuestas}

\subsection*{a. ¿Cuáles son los tres métodos que se utilizan para obtener información?}

Según el artículo de Kaspersky, los tres métodos son:

\begin{enumerate}[noitemsep]
    \item \textbf{Baiting:} se ofrece algo atractivo, como una memoria USB, para que la víctima caiga en una trampa.
    \item \textbf{Pretexting:} se inventa una historia o identidad para conseguir que la víctima entregue información.
    \item \textbf{Phishing:} se envían correos o mensajes falsos que imitan a una fuente confiable.
\end{enumerate}

\subsection*{b. Tres ejemplos de ataques de los dos primeros métodos}

\begin{itemize}[noitemsep]
    \item Dejar una memoria USB desconocida con malware para que alguien la conecte (baiting).
    \item Enviar una encuesta falsa que termina solicitando datos bancarios (pretexting).
    \item Hacerse pasar por un auditor o técnico para pedir información interna (pretexting).
\end{itemize}

\subsection*{c. ¿Las redes sociales son una amenaza? ¿Por qué?}

Sí. Pueden revelar nombres, familiares, lugares de trabajo, viajes y hábitos.
Un atacante puede usar esos datos para personalizar un engaño y parecer una
persona conocida. Por eso conviene limitar la información pública y revisar la
privacidad de las cuentas.

\subsection*{d. ¿Qué puede hacer una organización para defenderse?}

Debe capacitar periódicamente a sus empleados, establecer procedimientos para
verificar pedidos inusuales y usar contraseñas fuertes con autenticación
multifactor. También debe mantener actualizados sus sistemas, limitar los
permisos y realizar pruebas de phishing para detectar errores y corregirlos.

\subsection*{e. ¿Qué es el SANS Institute?}

El SANS Institute es una organización dedicada a la formación, investigación
y publicación de material sobre ciberseguridad. Ofrece cursos, certificaciones
y documentos técnicos para profesionales. El artículo consultado fue escrito
por Michael Alexander y se publicó en la biblioteca de documentos de SANS.

\end{document}
